
本文在统一储能与供需平衡模型下完成四个子问题的仿真。问题一的典型日购电量为$59482.70\ \mathrm{kWh}$，费用为$35126.95$元。问题二采用F1预测与风险裕量，334天总费用为$14786101.35$元，应急购电量为$245483.29\ \mathrm{kWh}$。问题三在0:00、6:00、12:00和18:00更新预报并调整合同，总费用为$14664345.40$元，相比不使用日内新预报的方案节省$566030.72$元，应急购电量为$210761.48\ \mathrm{kWh}$。问题四中，固定合同和可调合同的总费用分别为$17378687.69$元和$15465962.60$元，可调合同节省$1912725.09$元，并使应急购电量下降$63.14\%$。

上述结论仅适用于不可售电、无并网功率上限、充放电效率均为0.9、10分钟离散以及2025年历史样本的设定。数值检验显示供需平衡和SOC递推误差处于求解器浮点误差范围；敏感性分析表明，储能功率上限和预测误差倍率是影响费用的主要因素。模型还可扩展到含风电、需求响应和多储能的园区微网，但扩展后需要重新核验参数和费用口径。
